\section{Résultats de l'étape 2}

\subsection{Mode interactif}

L'étape 2 doit surtout être lue en mode interactif. Sur \texttt{data/galaxy\_1000}, avec \texttt{mpirun --bind-to none -np 2}, le temps total par image correspond ici à \texttt{Render time + MPI wait + point update time}.

\begin{table}[H]
    \centering
    \resizebox{\linewidth}{!}{%
    \begin{tabular}{c c c c c c}
        \toprule
        \textbf{Threads} & \textbf{Render (ms)} & \textbf{MPI wait + update (ms)} & \textbf{Total (ms)} & \textbf{Speed-up total} & \textbf{Efficacité totale} \\
        \midrule
        1 & 5.450 & 18.125 & 23.575 & 1.000 & 1.000 \\
        2 & 4.350 & 8.825 & 13.175 & 1.789 & 0.895 \\
        4 & 4.425 & 6.025 & 10.450 & 2.256 & 0.564 \\
        8 & 4.825 & 4.800 & 9.625 & 2.449 & 0.306 \\
        16 & 4.500 & 3.800 & 8.300 & 2.840 & 0.178 \\
        \bottomrule
    \end{tabular}%
    }
    \caption{Résultats du mode interactif pour l'étape 2}
    \label{tab:stage2_interactive}
\end{table}

On voit donc que le \texttt{rank 0} passe une partie du temps à attendre la nouvelle trame calculée par le \texttt{rank 1}, mais que le temps total reste du même ordre que celui de l'étape 1. La séparation MPI ne dégrade donc pas brutalement l'usage interactif.

\subsection{Benchmark MPI hors affichage}

Le tableau suivant ne représente pas le mode interactif. Il s'agit d'un benchmark MPI sans fenêtre, utilisé seulement pour estimer proprement le coût du calcul du \texttt{rank 1} et le surcoût de communication.

\begin{table}[H]
    \centering
    \resizebox{\linewidth}{!}{%
    \begin{tabular}{c c c c c c}
        \toprule
        \textbf{Threads} & \textbf{rank1\_compute (ms)} & \textbf{end-to-end (ms)} & \textbf{Surcoût MPI (ms)} & \textbf{Speed-up} & \textbf{Efficacité} \\
        \midrule
        1 & 246.802 & 246.919 & 0.117 & 1.000 & 1.000 \\
        2 & 138.732 & 138.853 & 0.121 & 1.778 & 0.889 \\
        4 & 79.649 & 79.767 & 0.118 & 3.096 & 0.774 \\
        8 & 52.546 & 52.673 & 0.127 & 4.688 & 0.586 \\
        16 & 42.668 & 42.801 & 0.132 & 5.769 & 0.361 \\
        \bottomrule
    \end{tabular}%
    }
    \caption{Résultats du benchmark MPI sans affichage}
    \label{tab:stage2_results}
\end{table}

\subsection{Lecture directe du benchmark}

Deux points ressortent immédiatement :
\begin{itemize}
    \item le temps \texttt{end-to-end} suit presque exactement le temps de calcul du \texttt{rank 1} ;
    \item le surcoût moyen de communication reste très faible, proche de \SI{0.12}{ms} par pas sur toute la campagne.
\end{itemize}

Autrement dit, la séparation MPI est fonctionnelle et son coût de transport reste négligeable à cette échelle. Le temps total du benchmark reste donc dominé par le calcul gravitationnel lui-même.

\subsection{Graphiques du benchmark}

\begin{figure}[H]
    \centering
    \begin{subfigure}[t]{0.48\linewidth}
        \centering
        \begin{tikzpicture}
\begin{axis}[
    width=\linewidth,
    height=6.1cm,
    xlabel={Nombre de threads sur le rank 1},
    ylabel={Speed-up},
    xmin=1,
    xmax=16,
    ymin=0,
    xtick={1,2,4,8,16},
    grid=both,
    minor y tick num=1,
    axis line style={draw=black!60},
]
\addplot[
    color=enstaBleuFonce,
    mark=*,
    line width=1.1pt,
]
coordinates {
        (1, 1.000)
        (2, 1.778)
        (4, 3.096)
        (8, 4.688)
        (16, 5.769)
};
\end{axis}
\end{tikzpicture}

        \caption{Speed-up bout-en-bout}
        \label{fig:stage2_speedup}
    \end{subfigure}
    \hfill
    \begin{subfigure}[t]{0.48\linewidth}
        \centering
        \begin{tikzpicture}
\begin{axis}[
    width=\linewidth,
    height=6.1cm,
    xlabel={Nombre de threads sur le rank 1},
    ylabel={Efficacité},
    xmin=1,
    xmax=16,
    ymin=0,
    ymax=1.05,
    xtick={1,2,4,8,16},
    grid=both,
    minor y tick num=1,
    axis line style={draw=black!60},
]
\addplot[
    color=enstaBleuClair,
    mark=square*,
    line width=1.1pt,
]
coordinates {
        (1, 1.000)
        (2, 0.889)
        (4, 0.774)
        (8, 0.586)
        (16, 0.361)
};
\end{axis}
\end{tikzpicture}

        \caption{Efficacité bout-en-bout}
        \label{fig:stage2_efficiency}
    \end{subfigure}
    \caption{Scalabilité observée sur le benchmark MPI hors affichage de l'étape 2}
    \label{fig:stage2_scaling}
\end{figure}
